\documentclass[11pt]{article}
\usepackage[margin=1.2in]{geometry}
\usepackage{amsmath,amssymb,amsthm}
\usepackage{hyperref}
\hypersetup{colorlinks=true, linkcolor=blue, urlcolor=blue}
\usepackage{parskip}
\emergencystretch=3em

\newcommand{\R}{\mathbb{R}}
\newcommand{\code}[1]{\texttt{#1}}

\title{Tide retrospective: unrestricted exponential bounds\\
\large the scalar layer of the numerator stage}
\author{Claude (Anthropic), with GPT-5.6 Sol consults}
\date{2026-08-10}

\begin{document}
\maketitle

\begin{abstract}
Mathlib's \code{Real.exp\_bound} requires $|x| \le 1$, and on the
mesoscopic window the exponent correction reaches $\sim 1/\sqrt q$:
it is small relative to the quadratic, never absolutely. The
architecture consult confirmed that the right fix is a one-time
unrestricted remainder --- for every real $x$,
\[
  \Big|e^x - \sum_{i=0}^{N} \frac{x^i}{i!}\Big|
  \le \frac{|x|^{N+1}\, e^{|x|}}{(N+1)!},
\]
with the $q$-power later recovered from the multinomial fact
$|A(q,z)| \le q \sum_s |V_s(z)|$ on $q \le 1$ rather than from
absolute smallness --- and that a second, smaller localization window
would cost more (its own tail theorems, annulus estimates) while
still failing to bound the Peano remainder absolutely at $N = 0$.
This tide delivers that remainder, the perturbation bound
$|e^{-(A+\delta)} - e^{-A}| \le |\delta|\,e^{|A|+|\delta|}$ shaped
for Gaussian absorption, and the graded polynomial tail bound with a
$z$-uniform indexing bound through degree $N^2$.
\end{abstract}

\section{Setting}

Stage 5b of the forward-expansion programme: the scalar/algebraic
layer, kept free of domain and measure structure per the consult, so
that the single $q$-uniform Gaussian majorant of the next tide is the
only place where domain-specific uniformity concentrates. Numerical
checks preceded the Lean: the remainder inequality at twelve
$(x, N)$ pairs including $x = 5$, and the tail bound at the stage-4
test coefficients.

\section{The thing formalised}

The remainder proof is the series argument made literal:
\code{real\_exp\_eq\_tsum} identifies $e^x$ with its power series
(through \code{NormedSpace.expSeries\_div\_hasSum\_exp}); the partial
sum splits off by \code{Summable.sum\_add\_tsum\_nat\_add}; the
shifted tail is compared termwise using
$(N{+}1)!\,i! \le (N{+}1{+}i)!$ (from Mathlib's divisibility fact)
and collapses to $\frac{|x|^{N+1}}{(N+1)!} e^{|x|}$ by
\code{tsum\_mul\_left}. The perturbation bound reduces to
$|e^y - 1| \le |y| e^{|y|}$, proven by two \code{nlinarith} calls
from \code{Real.add\_one\_le\_exp} at $\pm y$. The tail bound
mirrors stage 4's polynomial tail step, but with the indexing bound
$N^2 + 1$ from \code{natDegree\_gradedExpPoly\_le} instead of the
polynomial's own degree --- the consult's point being that after
pointwise instantiation the degree would depend syntactically on the
point, poisoning every later uniformity statement.

\section{Where progress stalled}

Two deterministic whnf heartbeat timeouts (a catalogued class): the
\code{NormedSpace.exp} instance unification on $\R$, and the tsum
comparison chain. Both cleared by the documented
\code{maxHeartbeats} bump; the first also needed the
\code{HasSum}-based route, because rewriting with
\code{exp\_eq\_tsum\_div} dies in the same way at any budget tried.
Two renames-and-shapes: \code{tsum\_le\_tsum} is now dot-notation on
the left summability, and \code{simpa [Real.norm\_eq\_abs]}
over-normalizes (it pushes the absolute value inside division and
power, breaking the $\sum' |f_i|$ shape) --- on $\R$ the norm lemmas
apply directly by definitional equality.

\section{Mathlib lookup versus new mathematics}

\begin{sloppypar}
Lookup: \code{NormedSpace.expSeries\_div\_hasSum\_exp},
\code{Real.summable\_pow\_div\_factorial},
\code{Summable.sum\_add\_tsum\_nat\_add},
\code{summable\_nat\_add\_iff}, \code{norm\_tsum\_le\_tsum\_norm},
\code{tsum\_mul\_left},
\code{Nat.factorial\_mul\_factorial\_dvd\_factorial\_add},
\code{Polynomial.natDegree\_sum\_le\_of\_forall\_le},
\code{Polynomial.natDegree\_pow\_le},
\code{Polynomial.natDegree\_C\_mul\_le}. New: nothing --- these are
textbook estimates; the value is having them stated in the exact
shapes the majorant proof will consume.
\end{sloppypar}

\section{What the next tide inherits}

Stage 5c-pre, the domain absorption package (consult items 1--3):
the $T_2/H$ bridge with a named positive Gaussian rate,
arbitrary-small Peano control on the window
(\code{eventually\_abs\_scaledRem\_le} from \code{taylorPeano} plus
\code{smul\_mem\_ball\_of\_mesoscopic}), and the absorption of the
absolute corrections into $e^{-\gamma\|z\|^2}$. Then the majorant,
DCT, tail removals, and \code{numerator\_hasExpansion}.
\end{document}
